\documentclass[12pt,a4paper]{article}

% ==========================
% ⚙️ 필수 패키지
% ==========================
\usepackage{fontspec}   % 폰트 설정
\usepackage{kotex}      % 한글 지원

\usepackage{setspace}   % 줄간격
\usepackage{titlesec}   % 제목 스타일
\usepackage{fancyhdr}   % 머리말/꼬리말
\usepackage{csquotes}   % 따옴표 처리
\usepackage{graphicx}   % 이미지
\usepackage{float}      % 그림/표 위치 고정
\usepackage{lscape}     % 가로 페이지
\usepackage{enumitem}   % 리스트 스타일
\usepackage[backend=biber,
            style=verbose-note,   % 각주형 노트 스타일
            autocite=footnote,    % \autocite 를 각주로
            ibidtracker=context,  % ibid 처리 선택사항
            doi=false,            % DOI 출력 비활성화
            sorting=nyt]{biblatex}
\addbibresource{references.bib} % BibTeX 파일
\usepackage{array, multirow, makecell, booktabs, longtable, tabularx} % 표 관련
\usepackage[table]{xcolor}
\newcolumntype{L}[1]{>{\raggedright\arraybackslash}p{#1}}
\newcolumntype{C}[1]{>{\centering\arraybackslash}p{#1}}
\newcolumntype{R}[1]{>{\raggedleft\arraybackslash}p{#1}}

\usepackage[most]{tcolorbox}  % 박스
\usepackage{algorithm, algorithmic} % 알고리즘
\usepackage{amsmath, amssymb, amsthm, gensymb, textcomp} % 수학 기호
\usepackage[babel]{microtype} % 자간
\usepackage{comment}          % 주석
\usepackage{subfigure}        % 서브 그림

% ==========================
% ✍️ 본문 서식
% ==========================
\setstretch{1.5}
\setlength{\parskip}{0.8em}
\setlength{\parindent}{15pt}
\setmainfont{Times New Roman}

\setmainhangulfont[
  Path = ./,
  UprightFont = *Batang Light.ttf,
  BoldFont    = *Batang Medium.ttf
]{KoPubWorld }

\setsanshangulfont[
  Path = ./,
  UprightFont = *Dotum Light.ttf,
  BoldFont    = *Dotum Medium.ttf
]{KoPubWorld }

\newfontfamily\hangulstrongfont[
  Path = ./,
  UprightFont = *Batang Bold.ttf
]{KoPubWorld }

% ==========================
% 📑 번호 및 목차
% ==========================
\renewcommand\thesection{\Roman{section}.}
% \renewcommand\thesection{\arabic{section}.}
\renewcommand\thesubsection{\arabic{subsection}.}
\renewcommand\theparagraph{\arabic{paragraph})}
\setcounter{secnumdepth}{1}

\setcounter{tocdepth}{3}

% 소제목 번호 (circled number)
\ExplSyntaxOn
\cs_new:Npn \circnum #1 { \char_generate:nn { "2460 + #1 - 1 } { 12 } }
\renewcommand\thesubparagraph{\circnum{\value{subparagraph}}}
\ExplSyntaxOff

\newcommand{\cl}[1]{\raisebox{.3pt}{\textcircled{\raisebox{-.9pt}{\small #1}}}}

% 알고리즘 input/output 라벨
\renewcommand{\algorithmicrequire}{\makebox[40px]{\hfill\textbf{Input :}}}
\renewcommand{\algorithmicensure}{\makebox[40px]{\hfill\textbf{Output :}}}

% ==========================
% 📄 목차 스타일
% ==========================
\usepackage{tocloft}
\renewcommand{\cfttoctitlefont}{\hangulstrongfont\Large}
\renewcommand{\cftaftertoctitle}{\vskip 1em}
\renewcommand{\cftsecfont}{\normalsize}
\renewcommand{\cftsubsecfont}{\small}
\renewcommand{\cftsubsubsecfont}{\small}
\setlength{\cftbeforesecskip}{2pt}

% ==========================
% 📄 제목 정보
% ==========================
\title{%
  \hangulstrongfont 법적 행위와 법적 행위자: 해석적 수렴, 제도적 역할, 그리고 법적 실천의 구조 \\[4pt]
  \normalsize (Legal Acts and Legal Agents: Interpretive Convergence, Institutional Roles, and the Structure of Legal Practice)
}

% ==========================
% 📄 페이지 스타일
% ==========================
\pagestyle{fancy}
\fancyhf{}              % 헤더/푸터 초기화
\fancyfoot[C]{\thepage} % 가운데 하단에 페이지만 표시
\renewcommand{\headrulewidth}{0pt} % 상단 라인도 제거

% ==========================
% 📚 본문 시작
% ==========================

\begin{document}
\maketitle
\tableofcontents  % 전체 목차 삽입
\newpage          % 목차 뒤에 새 페이지 시작

\begin{abstract}

본 논문은 해석적 수렴 접근을 법철학에 적용하는 작업의 두 번째 단계를
발전시키는 것으로, 법의 존재조건에서 \textbf{법적 실천의 구조} 및 다양한
법적 행위자의 역할로 초점을 이동한다. 이전 연구가 법적 규범성을
\textbf{해석적 수렴}의 산물로 이해한 데 기초하여, 본 논문은 개인의
구체적 \textbf{법적 행위} 및 제도적 \textbf{직무들}이 그러한 수렴에
어떻게 기여하는지, 그리고 사법부의 역할이 어떻게 체계적으로 오해되어
왔는지를 분석한다.

첫 번째 부분은 \textbf{법적 행위}의 일반 이론을 제시한다. 법적 행위란
그것을 통해 법의 성립을 지향하는 의사소통적 작용으로서, 다만 세 가지
적정성 제약(felicity constraints) --- \textbf{일반성},
\textbf{공유가능성}, \textbf{초점점 안정성} --- 의 충족을 조건으로 한다.
이러한 조건은, 입법 행위, 사법적 판결문, 행정적 조치, 준수(compliance)의
관행 중 어떤 것들은 법적 규범으로 안정화되는 반면, 어떤 것들은 주변적
위치에 머물거나 소멸하는 이유에 대한 설명을 제공한다. 두 번째 부분은
주요 \textbf{법적 행위자} --- 입법부, 행정부 공무담당자, 법원, 그리고
법의 수범자들 또는 시민들 --- 를 모두 해석적 수렴에 기여하는 법형성
과정의 참여자로 재개념화한다. 이를 통해 ``입법은 입법부의 기능, 해석은
법원의 기능''이라는 정설적 구도를 해체한다.

논문의 핵심은 \textbf{사법적 역할}을 해명하는데 있다. 본 논문은 법원이
법의 최종적 의미를 계시하는 형이상학적 신탁자(oracle)가 아니라, 그
결정이 법적 공동체 내에서의 후속적 화행적수인(uptake)을 통해서만 더 넓은
권위를 획득하는 \textbf{분쟁 해결 제도(dispute-resolution
institutions)}로 이해되어야 한다고 주장한다. 사법적 우위, 난해한 사안,
그리고 소위 ``사법적극주의(activism)''를 둘러싼 익숙한 난제들은, 해당
판결이 기존의 전문적으로 제도화된 규범, 선례 존중 관행, 사법에서의 심사
구조 속에서 G--S--F 조건을 얼마나 충족하는지에 관한 질문으로 재구성된다.
마지막 부분은 이러한 분석이 권력분립, 선례의 규율, 그리고 사법적
정당성(judicial legitimacy)에 전달하는 함의를 도출한다. 본 논문은 해석적
수렴 이론의 기초를 반복하기보다, 그러한 틀이 법적 행위, 법적 행위자,
그리고 법적 실천의 내부 구조에 대한 이해를 어떻게 새롭게 구성하는지를
보여주는 데 목적을 둔다.

\end{abstract}

\section{서론}

어떻게 수많은 독립적 행위자들이 여러 제도에 걸쳐 존재하는 법체계가
그럼에도 불구하고 마치 하나의 일관된 목소리로 말하는 것처럼 기능할 수
있는가? 이 질문은 법에서의 \textbf{해석적 수렴}에 관한 근본적 난제를
제기한다. 판사, 입법자, 행정담당자, 시민 등 각각이 법적 규범의 창출과
적용에 기여함에도 불구하고, 그들이 법의 \textsf{의미(meaning)}와
\textsf{힘(force)}에 수렴할 것이라는 기대가 존재한다. 다시 말해, 법적
실천은 서로 다른 \textbf{법적 행위자(legal agents)}들이 \textbf{법적
행위(legal acts)}를 상호 이해 가능하고 일관된 방식으로 파악하고 수행할
수 있음을 전제한다. 그러나 해석이 분기할(diverge) 가능성은 매우 크며,
특히 \textbf{사법적 행위자}의 경우에서 두드러진다. 이들은 난해한
사안(hard cases)에서 해석의 자유를 가지며, 이는 법의 지배(rule of law)가
요구하는 바로 그 수렴을 위협하는 듯 보인다.

본 논문의 목적은 이 난제를 재구성하고 해결하는 데 있다. 이를 통해 우리는
\textbf{법의 해석적 수렴 이론}을 기초부터 정교화하며, 그 기저에 놓인
문제를 해결하는 과정을 드러내고자 한다. 이 이론의 핵심 생각은, 법적
행위의 의미는 어느 한 행위자의 사적 발명이 아니라 해석자들의 공동체
내에서의 공유된 해석적 실천과 화행적수인(uptake)의 산물이라는
점이다.\footnote{`화행적수인(uptake)'은 발화행위 이론에서 발화수반행위
  (illocutionary act)의 성립과 밀접하게 연관된 개념으로, 존 오스틴(John
  Austin)이 \emph{How to Do Things with Words}에서 발화수반행위의 수행이
  ``화행적수인의 확보(securing uptake)''를 필수적으로 요구한다고 설명한
  데 그 중요한 이론적 기원을 둔다\autocite{Austin1962}. 이후 P.F.
  스트로슨(P.F. Strawson)은 이를 의도와 관습의 관계 속에서
  재논의하였으며\autocite{Strawson1964}, 마리나 즈비자(Marina Sbisà)는
  화행적수인의 불가결성을 발화수반력(illocution)의 표지로 재해석하면서
  오스틴, 스트로슨, 서얼 등의 독해를 비판적으로
  정리하였다\autocite{Sbisa2009}. 최근에는 화행적수인 자체가 독립된 문제로
  다시 조명되면서, 오스틴 이후의 논쟁을 재구성하려는 논의도 전개되고
  있다\autocite{deLara2024}.} 도전 과제는 이것이 구체적으로 어떻게
작동하는지를 명확히 하는 것이다. 다양한 법적 행위자들은 어떻게 해석을
조정하는가(coordinate)? 어떤 제도적 장치가 수렴을 보장하며, 취약점은
어디에 있는가? 무엇보다도, 가장 두드러진 해석적 역할이라 할 수 있는
사법적 행위자의 역할을 어떻게 이해해야 흔히 제기되는 오해들을 해소할 수
있는가?

논의는 다음과 같은 순서로 진행된다. \textbf{제II절}에서는 법적 행위와
법적 행위자의 개념을 정의함으로써 이론의 기초를 마련하고, 실천-매개적
해석 상호작용을 모델링하는 \textbf{G--S--F 적정성 제약(felicity
constraints)}과 \textbf{화행적수인 함수(uptake functions)}를 도입한다.
\textbf{제III절}에서는 중심 난제를 제시한다. 즉, 법적 실천에 해석적
수렴이 필수적이라는 점과 해석의 분기가 명백히 실재한다는 점 사이의
긴장을, 특히 역설적으로 보이는 사법적 역할에 초점을 맞추어 검토한다.
예컨대, 제정법을 단순히 입법부에서 사법부로 전달되는 메시지(message)로
보는 순진한 의사소통관은 왜 실패하는지, 그리고 이것이 왜 곤경에
빠지는지를 설명할 것이다. \textbf{제IV절}에서는 단계적 해결을 제시한다.
우리는 언어철학의 통찰---특히 규칙 따르기에 관한 루드비히
비트겐슈타인(Ludwig Wittgenstein)의 논의와 발화행위(speech act) 이론에
관한 J. L. 오스틴(J. L. Austin)의 논의---을 활용하여, 법적 의미는 고립된
의지적 행위의 산물이 아니라 공유된 해석적 실천과 구조화된 화행적수인의
산물임을 논증한다. 특히, 우리는 사법적 역할을 명확히 한다. 판사는 의미의
무제약적 원천도 아니고 단지 입법부의 뜻을 전달하는 복화술사도 아니다.
그는, 역할에 의해 구조화되고 실천을 구조화하는, 수렴을 촉진하는 제약
속에서 활동한다.

우리는 법에서의 의사소통적 내용의 성격, 화자 의미(speaker meaning)와
공동체 규범의 상호작용, 다수 해석자 간의 간극을 메우는 제도적 설계와
같은 핵심 쟁점들에 직접적으로 맞선다. 각 단계에서 하나의 테제나 잠정적
해결책을 제시하고, 그 결점을 검토하며, 이에 응답하여 우리의 이해를
명료히 할 것이다. 이를 통해 우리는 해석적 수렴이 어떻게 가능한지, 그리고
사법적 역할에 대한 오해가 어떻게 그 그림을 가려왔는지를 보다 분명히
보여주고자 한다. 결론에 이르면, 의견이 통합되지 않는 행위자들로부터
통일된 법적 실천이 출현하는 당혹스러운 현상의 신비함은 훨씬 누그러져
있을 것이다.

\section{기초: 법적 행위, 행위자, 그리고 해석적 수렴
이론}

\subsection{1. 법적 행위}

\textbf{법적 행위(Legal acts)}란 그것을 통해 법이 \textsf{성립하게
되는(comes to be)} 의사소통적 작용을 말한다. 즉, 해석자들의 공동체
내에서 화행적으로 수인되고 안정화될 때, 법의 내용과 권위를
\textsf{구성하게 되는} 행위들이다.\footnote{이 설명은 법에 대한 관행주의적
  접근 및 조정 기반 접근에 일정 부분 빚지고 있지만, 그것들과 동일한 것은
  아니다. 루이스(Lewis)는 반복되는 조정 문제 속에서 관행을 자기유지적인
  정기성(regularity)으로 분석하고\autocite{Lewis1969}, 포스티머(Postema)는
  법에 대한 조정과 관행 설명을 전개한다\autocite{Postema1982}. 이에 비해 본
  이론은 관행이나 균형 그 자체보다, 법적 행위의 적정성과 법적 행위자들의
  분화된 화행적 수인으로 초점을 이동한다.} 실천 속에서 법적 의미를
안정화하는 데 기여하는 행위들의 대표적인 예로는 입법적 제정, 사법적
결정, 행정적 조치, 그리고 수범자들의 준수나 의존 패턴 등이 있을 수 있다.
이러한 각 행위는 의무와 권리를 생성·변경·적용함으로써 규범적 지형을
변화시키는 것을 의도한다. 핵심적으로, 법적 행위는 제도적 지위에 의해
\textsf{사전에} 정의되는 것이 아니라, 특정한 \textbf{적정성 조건(felicity
conditions)} --- 즉, 한 행위가 해석공동체 내에서 법으로서
화행적수인되고, 공유되며, 안정화될 수 있는 조건들 --- 을 충족함으로써
\textsf{제도화 되어간다(becomes institutional)}. 우리는 이러한 조건들을
\textbf{G--S--F}라고 명명하기로 한다. G, S, F의 각 영문 이니셜은 다음의
\textbf{일반성(Generality)}, \textbf{공유가능성(Shareability)},
\textbf{초점점 안정성(Focal-point Stability)} 을 나타낸다.\footnote{G--S--F
  세 조건은 여기에서 독창적인 재구성으로 제시되지만, 그것은
  법다움(legality)이나 법의 지배(rule of law) 이론에서 익숙한 한 계보와
  인접해 있다. 론 풀러(Lon Fuller)는 법다움의 중심 특징으로 일반성,
  공지성(publicity), 명료성을 전면에 내세우고, 조셉 라즈(Joseph Raz)는
  법의 지배가 일반적이고, 공개되어 있으며, 명확하고, 비교적 안정적인
  규칙을 요구한다고 강조한다. 본 이론은 이러한 생각을 법체계 전체의
  수준에서 다시 진술하는 것이 아니라, 그것들을 법적 행위와 법적 실천
  내부에서의 화행적수인에 관한 적정성 제약으로 재구성한다.
  \autocite{Fuller1969,Raz1979RuleOfLaw}}

\begin{itemize}
\item
  \textbf{일반성(Generality, G)} -- 행위의 내용은 단순히 일회적인 지시가
  아니라 \textbf{규칙, 원리, 또는 그 밖의 이후에 투사 가능한 규범적
  패턴(projectable normative pattern)}의 성격을 가져야 한다. 그것은 여러
  사례에 걸쳐 확장될 수 있어야 하며, 그래야 다른 행위자들이 그것을
  미래로 투사하여 행동, 정당화, 또는 의존의 이유로 삼을 수 있다. 이러한
  규칙적이거나 그 밖의 일반화 가능한 형식이 없다면, 어떠한 규범적
  화행적수인도 안정화될 수 없다.\footnote{이 요건은 풀러의 일반성 원리
    및 법이 일반 규칙을 통해 인도해야 한다는 라즈의 주장과 분명한
    친연성이 있다. 그러나 여기서의 주장은 법적 행위의 수준에서 다시
    구성된다. 즉 ``일반성''이란 규칙, 원리, 또는 판결이유와 같은
    형식으로 그 행위의 내용이 이후의 해석과 행위로 투사될 수 있는
    가능성을 가리킨다. \autocite{Fuller1969,Raz1979RuleOfLaw}}
\item
  \textbf{공유가능성(Shareability, S)} -- 제안된 규범은 \textbf{공적으로
  이해 가능하고 공유된 법적 어휘로 표현 가능해야} 하며, 그래야 참여자들
  사이에서 논박되고, 옹호되며, 정당화될 수 있다. 공유가능성은 이유와
  의미가 사적 의도에 갇힌 채 남아 있는 것이 아니라 공동체의 공동
  재산(common property)이 된다는 것을 뜻한다. 공유 가능해야 하는 것은
  텍스트만이 아니라, 그것을 이해하고 적용하기 위해 제시되는 이유들까지
  포함한다.\footnote{공유가능성은 공지성(publicity)이나
    공포(promulgation)와 겹치지만, 그것으로 다 소진되지는 않는다. 풀러는
    공지성과 이해 가능성을 법다움의 핵심으로 다루고, 라즈 역시 법은
    공개되어 있고 명확해야 한다고 강조하며, 제러미 월드론(Jeremy
    Waldron)은 법적 이유가 진술되고, 논박되며, 응답되어야 하는
    공적·절차적 차원을 부각한다. 본문의 조건은 공유 가능해야 하는 것이
    단지 텍스트만이 아니라, 공유된 법적 어휘 속에서 그 텍스트를 위해
    제시되는 이유들까지 포함한다는 점을 추가한다. \autocite{Fuller1969,Raz1979RuleOfLaw,Waldron2011}}
\item
  \textbf{초점점 안정성(Focal-point Stability, F)} -- 그 행위의 내용은
  시간이 경과해도 기대와 행동의 \textbf{조정 초점점(coordination focal
  point)}으로 기능할 수 있을 만큼 \textbf{충분히 안정적이고
  현저해야(stable and salient enough)} 한다. 어떤 초점적 해석이 형성되고
  실천을 통해 유지될 때, 그 규범은 자기강화적인 성격을 띠게
  된다.\footnote{이 조건의 초점점 요소는 현저성, 관행, 그리고 법적
    조정에 관한 보다 넓은 문헌에 기대고 있다. 셸링(Schelling)의 고전적
    설명은 상호의존적 상황에서 현저한 해결책(salient solutions)이 어떻게
    조정을 가능하게 하는지를 보여 주고\autocite{Schelling1960}, 루이스는
    반복되는 조정 문제 속에서 관행을 자기유지적인 정기성으로 보는 관련된
    구상을 발전시킨다\autocite{Lewis1969}. 법이론에서 포스티머(Postema)는
    법을 조정과 관행에 연결시키고\autocite{Postema1982}, 맥애덤스(McAdams)는
    법이 하나의 결과를 현저하게 만듦으로써 표현적(expressive)으로 작동할
    수 있다고 주장하며\autocite{McAdams2000}, 맥애덤스와 내들러(McAdams and
    Nadler)는 제재(sanctions)가 없더라도 법규칙이 준수를 위한 초점점을
    형성할 수 있다는 주장에 실험적 뒷받침을
    제공한다\autocite{McAdamsNadler2005,McAdamsNadler2008}. 또한
    비스트라노프스키(Bystranowski), 한니카이넨(Hannikainen), 그리고
    토비아(Tobia)는 조정 분석을 법해석으로
    확장한다\autocite{BystranowskiHannikainenTobia2025}. 본문의 주장은
    이러한 설명들에 기반하고 있지만, 그것들로 환원되지는 않는다. 즉 ``법
    준수(legal compliance)에 관한'' 초점점 이론과 달리, 본 논문은 초점점
    안정성을 일반적인 법적 행위에 관한 하나의 적정성 제약으로 다룬다. 그
    안정성 요소는 비교적 안정적인 규칙을 강조하는 법의 지배에 대한
    설명들과도 친연성을 가지지만, 여기서는 그러한 생각이 법체계 전체의
    수준이 아니라 행위-기반 수인의 수준에서 재구성된다.
    \autocite{Raz1979RuleOfLaw}}
\end{itemize}

이러한 조건들은 사전에 존재하는 제도적 질서(pre-existing institutional
order)를 전제하지 않는다. 전통적으로 법적 행위는 승인된
\textbf{역할}(입법, 사법, 행정 등)과 확립된 절차적 요건 아래에서
수행되는 제도적 행위로 이해되어 왔다. 존 서얼(John Searle)의 잘 알려진
구성 규칙(constitutive rules) 정식화---``X가 맥락 C에서 Y로 간주된다(X
counts as Y in context C)''\autocite{Searle1995}---는 이러한 제도적 현상을
설명하는 전통적인 방식을 잘 보여준다. 이러한 설명에 따르면, 법체계는
예컨대 다음과 같은 규칙들에 의해 구조화된다. 즉 \textsf{권한 있는 입법부가
특정한 단어를 발화하는 행위(X)가 법적 맥락(C)에서 법률 규칙의 창출(Y)로
간주된다}거나, \textsf{적법하게 임명된 판사가 판결을 선고하는 행위(X)가
구속력 있는 법원 명령(Y)으로 간주된다}는 식이다. 이러한 의미에서 구성
규칙은 특정한 수행이 법적으로 중요한 행위로 인정되는 방식을 설명하며,
이미 확립된 제도적 틀 안에서 어떤 행위들이 법적 의미를 갖는지
규정한다.\autocite{Searle1969}

성숙한 법체계에서는 실제로 많은---어쩌면 대부분의---법적 행위들이 이러한
방식으로 기술될 수 있다. 입법 제정, 사법적 판결, 행정적 명령과 같은
행위들은 누가 입법할 수 있는지, 판결이 어떻게 선고되는지, 무엇이
공포(promulgation)로 간주되는지 등을 규정하는 구성적 관행에 의해
안내되는 제도적 수행행위(institutional performatives)처럼 보인다. 이러한
관점에서 보면, 법적 행위의 효력은 어떤 행위들이 법적으로 유효한 수행으로
간주되는지를 정의하는 제도적 규칙들로부터 비롯되는 것처럼 보인다.

그러나 이러한 설명은 법적 실천의 기초적 발생 논리를 제시한다기보다는,
이미 안정화된 법적 실천을 사후적으로 기술하는 방식으로 이해되어야 한다.
G--S--F 조건은 행위들이 해석적 수렴을 통해 어떻게 제도적 효력을 획득하게
되는지를 설명한다. 이러한 점에서 그것들은 서얼의 의미에서의 제도적 행위
설명과 유비되기보다는, 오히려 오스틴이 일상 언어에서 제시한
\textbf{발화행위의 적정성 조건}과 더 가까운 유비를 이룬다. 오스틴은
약속하기나 명령하기와 같은 발화행위를 수행하기 위해서는 단순히 특정한
단어를 발화하는 것만으로는 충분하지 않으며, 적절한 여건, 화자의 권위,
그리고 청중의 화행적수인(audience uptake)이 갖추어져야 그 행위가 성공할
수 있음을 보여주었다.\autocite{Austin1962}

이와 유사하게 \textbf{G--S--F} 제약은 해석 공동체 내부에서 행위가
\textbf{화행적 수인되고 안정화될 수 있는 적합성}을 가리킨다. 인정된
절차적 관행들(예컨대 적절한 법정, 정족수, 형식)을 충족하는 것은 행위가
실제로 작동하는 규범이 되기 위한 \textbf{필요조건일 수는 있지만
충분조건은 아니다}. 그 내용이 특정한 사건을 넘어 일반적으로 적용될 수
있을 만큼 규칙적일 때(\textbf{일반성}), 그 이유가 공유된 법적 어휘
속에서 공적으로 이해되고 논증될 수 있을 때(\textbf{공유가능성}), 그리고
시간의 경과에 따라 안정적인 조정 초점으로 기능할 만큼 충분히 안정적일
때(\textbf{초점점 안정성})에만 비로소 그 행위는 법적 실천의 규범적
네트워크 속에 실제로 편입된다.

이 관점은 익숙한 실패 양상들도 설명해 준다. 형식적으로는 유효한 법률이라
하더라도 극도로 특정적인 사안에 대응하는 방식으로 제시된 경우에는 구성
검사는 통과하더라도 \textbf{일반성}을 충족하지 못한다. 판결이 효과적으로
전달되지 않거나 불분명한 관용구 등으로 표현된 경우에는
\textbf{공유가능성}을 충족하지 못한다. 특정 법리가 재판부가 바뀔 때마다
흔들리는 경우에는 \textbf{초점점 안정성}을 충족하지 못한다. 이러한 경우
행위는 제도적 형식 요건을 충족하더라도 실제 법적 실천 속에서 법으로
안정화되는 데에는 실패할 수 있다. 이러한 의미에서 인정된 권위와 절차
자체도 G--S--F가 설명하는 동일한 수렴 과정의 산물이며, 그것들은 법적
힘의 선행적 원천이라기보다는 성공적인 화행적수인의 패턴 속에서 출현한
안정화된 구조의 일부로 이해되어야 한다.

이들 중 어느 하나라도 실패하면, 그 행위는 불발되거나(misfire) 사회적으로
비활성적인 상태(socially inert)에 머무를 수 있다.\footnote{여기에서
  ``불발된다''는 표현은 관련된 수행적 절차가 적절하게 호출되거나
  실행되지 않아 그 행위가 ``무효이거나 효력이 없는 것(void or without
  effect)''이 되는 경우를 가리키는 오스틴의 기술적 의미로 사용하고 있다.
  \autocite{Austin1962} ``사회적으로 비활성적''이라는 표현은 
  형식적으로는 인식 가능한 행위일지라도 해석 공동체 안에서 법으로
  안정화되는 데 실패하는 경우를 설명하기 위해 이를 확장한 것이다.
  } 미공표 판결은 \textbf{공유가능성}의 결여로,
과도하게 특수한 부칙 조항은 \textbf{일반성}의 결여로 실패한다. 또한
법원마다 판시가 오락가락하거나 신뢰가 형성되기 전에 반복적으로 번복되는
법리는 \textbf{초점 안정성}의 결여로 실패한다. \textbf{권위}와
\textbf{절차}는 선행적 주어진 결정된 요소가 아니라, 이러한 수렴적 실천에
의해 \textbf{안정화된 효과} 그 자체이다. 따라서 G--S--F 세 조건은
행위들이 해석적 수렴을 통해 \textsf{법이 되어가는(become law)} 과정에서
작동하는 \textbf{내용 및 실천 지향적 적정성 제약}을 포착한다.

\subsection{2. 법적 행위자}

\textbf{법적 행위자(legal agents)}란, 그들의 수행, 응답, 그리고 다양한
형태의 화행적수인이 법의 형성과 안정화에 참여하는 행위를 하는 자들을
말한다. 이 설명에 따르면 법적 행위자라는 개념은 미리 고정된 제도적
질서를 선행적으로 전제하지 않는다. 그러나 성숙한 법적 실천에서는
행위자들이 보통 공적으로 인정된 직무에 기반한 역할 또는 실천 속에서
인정된 역할을 차지하며, 그에 상응하는 권한과 책임을 가진다. 현재의
논의를 위해 고려되는 주요 행위자들은 다음과 같다. (1) 입법기관 ---
일반적 규범(general norms)을 정식화한다. (2) 법원 --- 분쟁을 해결하고,
이후의 실천에서 초점점이 될 수 있는 이유들을 생성한다. (3) 행정부 및
행정 공무담당자 --- 법적 규범을 집행 및 실행하고, 이를 추가적으로
구체화한다. (4) 법의 수범자들 --- 시민, 기업, 결사체 및 기타 법의
수신자들을 포함하며, 그들의 준수, 이의제기, 계약 체결, 소송제기, 법에
대한 의존 패턴 등이 법적 의미를 안정화하거나 또는 불안정하게 만드는 데
기여한다. 이들의 기여 방식은 서로 다르지만, 법적 실천이 하나의 정합적인
체계로 유지되는 것은 오직 이러한 서로 다른 행위와 화행적수인의 형태들이
상호 응답하며 연결되어 있고, 시간의 경과에 따라 수렴적 정렬 속으로 다시
묶일 수 있을 때에 한정된다.

\subsection{3. 화행적수인 함수와
수렴}

우리는 이 이론의 핵심 측면을 \textbf{화행적수인 함수(uptake function)}
개념을 통해 형식화할 수 있다. 발화행위 이론에서 차용한
``화행적수인''이라는 용어를 이용하여, 우리는 해석을 일종의 \textsf{응답
함수(response function)}로 취급한다. 즉, 입력(이전의 법적 행위 또는
메시지)을 출력(해석자의 이해 또는 실행된 결정)으로 사상하는(mapping)
함수로 이해하는 것이다. 각 법적 행위자 \(A\)에 대하여, 화행적수인 함수
\(U_A\)를 정의하면 \(U_A(x)\)는 행위자 \(A\)가 다른 행위자가 산출한 법적
내용 \(x\)를 어떻게 해석하거나 사용하는지를 나타낸다. 예컨대, \(x\)가
입법부가 제정한 법률이라면, \(U_{\text{judge}}(x)\)는 특정 사건에서
판사(judge)가 그 법률을 어떻게 해석하는지를 의미한다.
\(U_{\text{police}}(x)\)는 행정 행위자로서의 경찰(police officer)이 그
법률 하에서 자신의 권한을 어떻게 이해하는지를 나타내며,
\(U_{\text{citizen}}(x)\)는 시민(citizen)이 그 법률에 따라 자신의 의무를
어떻게 이해하는지를 나타낸다. 이 이론의 테제는, 법적 실천이 작동하기
위해서는, 이러한 다양한 화행적수인 함수들이 상호 양립 가능한 출력들로
\textbf{수렴하여야} 한다는 것이다. 단순화하여 말하면, 관련된 화행적수인
함수들은 공유된 해석 기준 속에서 안정화된 법의 공적으로 접근 가능한
내용과 양립 가능한 산출로 수렴해야 한다. 판사, 행정담당자, 그리고 법적
주체들은 어떤 숨겨진 입법자의 심리를 재현할 필요는 없다. 오히려 그들의
화행적수인은 법 규범이 행위를 향도하고, 결정을 정당화하며, 기대를 조정할
수 있을 만큼 충분히 긴밀하게 정렬되어 있어야 한다. 만약 \(U_A\)와
\(U_B\)가 동일한 입력 \(x\)에 대해 서로 양립할 수 없는 결과를 반복적으로
산출한다면, 법은 일관성 없는 조각들로 분열될 것이다. 따라서 해석적 수렴
조건은 단일하고 영원히 고정된 균형을 가리키는 것이 아니라, 관련
행위자들의 수인 함수들 사이에 형성되는 충분히 안정적인 초점적 정렬을
가리킨다.

이러한 수렴은 자동적이지도, 사소하지도 않다. 이는 법적 실천의 성취로,
설명이 요구된다. 단순히 성문화된 제정법이나 선례가 존재한다고 해서
수렴이 보장되는 것은 아니다. 마치 언어(language)에 단어들이 존재한다고
해서 의사소통이 보장되지 않는 것과 같다. 입법부는 텍스트를 공포할 수
있지만, 그 텍스트가 다른 이들에 의해 어떻게 해석되는지는 실천과 관행의
문제이다. 실제로 상황은 공동체가 자연언어(natural language)를 사용하는
것과 유사하다. 데이비드 루이스(David Lewis)가 관찰했듯, 한 집단이 언어
\(L\)을 사용한다는 것은 \(L\) 안에서 진실성과 신뢰의 관행(convention of
truthfulness and trust)에 따르는 것이다. 즉, 화자는 일반적으로 공유된
언어에 따라 진실되게 말하려는 의도를 가지며, 청자는 화자가 그렇게 하고
있다고 신뢰한다\autocite{Lewis1969}. 이러한 관행은 조정 해법(coordination
solution)이다. 모든 사람은 이해되기 위해 단어를 일관되게 사용하는 데에
이해관계가 있다\autocite{Lewis1975}. 마찬가지로, 법적 행위자들은 법에 대한
충실성(fidelity to law)의 관행(명시적이든 묵시적이든)에 참여한다.
입법자는 자신의 제정이 특정한 방식으로 적용되기를 의도하며, 판사는
제정된 법을 합리적이고 제도적으로 승인된 방식으로 해석하려고 노력하고,
시민·공무담당자들은 법적 권위자들이 확립된 해석에서 갑자기 이탈하지 않을
것이라고 신뢰한다. 이러한 관점에서 보면, \textsf{해석적 수렴}은 법적
공동체 내에서 스스로 안정화하는 조정 평형상태로 이해될 수 있다. 즉, 법적
텍스트를 일관된 방식으로 사용하는 상호적 일치(mutual agreement)이며,
이는 모든 당사자가 공유된 이해의 이익(benefits)과 혼란의 비용(costs)을
인식하기 때문에 유지된다.

논의의 골격을 요약해 보면, 우리는 법적 행위와 행위자를 정의하고, 법적
행위의 적정성 조건을 포착하기 위해 G--S--F 조건을 도입했으며, 법적
내용이 행위자들 간에 어떻게 전달되는지를 모델링하기 위해 화행적수인
함수를 설정했다. 이러한 도구들을 갖추고, 이제 우리는 핵심적 난제로
나아간다. 즉, 왜 우리는 해석적 수렴이 유지된다고 기대해야 하는가? 그리고
그것이 붕괴되는 듯 보일 때는 어떻게 설명해야 하는가? 보다 구체적으로
말하자면, 해석적 수렴이 본질적이라면, \textbf{사법적 역할은 왜 그렇게
문제적으로 보이게 되었는가}? 다음 절에서는 이 도전을 상세히 정식화한다.

\section{해석적 수렴의 난제: 분기, 재량, 그리고 사법적
역할}

언뜻 보기에, 법의 통일된 해석을 보장하는 것은 단순히 모든 법적 행위자가
입법자의 의도에 복종하는 문제라고 생각할 수도 있다. 예컨대, 전통적으로
\textbf{충실한 대리인 모형(faithful-agent model)}은 판사를 입법부의
``충실한 대리인(faithful agents)''으로 간주하며, 판사는 법률(statutes)에
표현된 입법부의 의도를 수행하는 임무를 맡는다고
본다\autocite{bruhl2012elected,ross2013reconsidering}. 이러한 순진한
관점에서는 모든 이가 단일한 권위적 의미, 즉 입법부가 마음속에 가진
의미에 복종하기 때문에 수렴이 발생한다고 본다. 각 판사가 ``입법부는
여기서 무엇을 의도했는가?''라고 묻고, 각 시민이 ``법원이 이 법률을
고려할 때 무엇이라고 말할까?''라고 묻는다면, 모두가 입법부의 전달된
의도라는 동일한 답에 도달할 것이라는 가정이다.

그러나 더 면밀히 살펴보면, 이 그림에 균열이 생기기 시작한다.
\textbf{의도는 스스로 해석되지 않는다.} 입법 의도는 (그런 것이
존재한다고 가정하고 발견 가능하다 하더라도) 언어를 통해 전달되어야
하는데, 언어는 본질적으로 애매성(ambiguity)과 불명확성(vagueness)에
노출되어 있다. 입법부는 일반적으로 매우 포괄적인 용어로 말하며, 이는
입법자가 명시적으로 예상하지 못했던 수많은 상황에 적용된다. 이로 인해
하트(H. L. A. Hart)가 지적한 바와 같이, 법 규칙에는 ``언어의 개방적
직조(open texture)'' 때문에 \textbf{명확한 의미의 핵심부(core of clear
meaning)}와 \textbf{의심스러운 반영부(penumbra of doubt)}가
존재한다\autocite{hart1994concept}. 예를 들어, 전동 킥보드가 시 조례에서
말하는 ``공원의 탈것(vehicle)''에 해당하는지와 같은 난해한 사안이나 경계
사례(borderline cases)에서는, 선의의 해석자들조차 해석이 분기될 수 있다.
입법 의도가 유일한 해답을 산출하지 못하는 경우도 있는데, 이는 입법자들이
해당 사례를 고려하지 않았거나, 그들의 의도가 일관되지 않거나 불명확했기
때문이다.

더 나아가, 판사들이 단순히 입법 의도를 \textsf{발견하여(discover)}
적용한다고 보는 가정은 \textbf{사법적 재량(judicial discretion)}의
현실에 의해 복잡해진다. 하트는 법적 자료들이 사건의 결론을 충분히
결정하지 못할 경우, 판사는 불가피하게 재량을 행사해야 하며, 사실상 그
간극을 메우기 위해 \textbf{새로운 법}을 창출하게 된다고
지적했다\autocite{hart1994concept}. 이는 하트와 드워킨(Ronald Dworkin) 간의
핵심 쟁점이었다. 하트는 난해한 사안에서 판사가 일정 부분 법을 창출한다는
점을 수용한 반면, 드워킨은 난해한 사안에도 법 안에 내재된 정답이
존재한다고 주장하며 진정한 재량은 부정했다\autocite{dworkin1977rights}.
우리의 목적에서 중요한 점은, 판사가 때때로 단순한 해석을 넘어 실제로
혁신을 수행한다면, 수렴은 어떻게 유지될 수 있는가 하는 것이다. 서로 다른
판사들이 서로 다른 방향으로 혁신한다면 분기가 발생하지 않겠는가? 설령 한
판사의 혁신이 선례가 된다 하더라도, 입법부가 의도한 것과 판사가 법으로
공표한 것 사이의 연계는 희미해질 수 있다.

이 난제는 \textbf{다중 행위자(multi-agent) 측면}을 고려할 때 더욱
심화된다. 법은 단일한 화자가 단일한 청자에게 보내는 일회적 메시지가
아니라, 다수의 당사자들 사이에서 이루어지는 지속적인 의사소통의
연쇄(chain of communications)이다. 입법부는 법률을 통해 소통하고,
사법부는 판결을 통해, 행정가는 행정규칙과 명령을 통해 소통한다. 시민들은
법의 내용을 일부는 법률에서, 일부는 공무담당자들이 그것을 어떻게
집행하는지에서 배운다. 이 되풀이 과정(iterative process)은 마치 전달하기
게임(game of telephone)와 같으며, 각 단계마다 왜곡의 위험이 존재한다.
다만 법체계는 왜곡을 줄이기 위한 견제 장치를 갖추고 있다. 문제는, 그
장치들이 충분한가 하는 것이다. 중요한 행위자 ― 이를테면 판사 ― 가 법을
일탈적 방식으로 해석한다면, 이후의 행위자들(하급 법원, 공무담당자,
대중)은 그 일탈적 해석을 따르고 결국 입법부의 의미에서 집단적으로
이탈하거나, 또는 그 해석과 충돌하고 결국 수렴을 상실할 수 있다. 양쪽
결과 모두 문제가 된다. 전자는 법의 민주적 원천을 약화시키고, 후자는 법의
지배의 일관성을 훼손한다.

문제를 단도직입적으로 제기하자면 다음과 같다. \textbf{``법적 해석이
바벨탑의 혼란으로 빠지지 않도록 보장하는 것은 무엇인가?''} 우리는 해석적
수렴이 규범이자 이상이라고 가정하지만, 그것이 어떻게 유지되는지를
설명해야 한다. 언어철학자들은 오랫동안 언어적 의미에 관한 유사한 질문과
씨름해 왔다. 곧, 현재의 어떤 사용도 무수히 다른 해석들과 양립 가능할 수
있는데, 그럼에도 우리는 어떻게 모두가 동일한 단어 사용 규칙을 따를 수
있는가? 잘 알려져있듯, 비트겐슈타인의 규칙 따르기 역설은 규칙이 그
자체로 올바른 적용에 관한 지침을 내포한다는 관념에
도전한다\autocite{wittgenstein1953,kripke1982}. 사울 크립키(Saul Kripke)는
비트겐슈타인을 해석하면서, 규칙을 따를 때 다음에 무엇을 해야 하는지를
결정해 줄 수 있는 사적이고 정신에 국한된 사실(private, mind-bound
fact)은 존재하지 않으며, 오직 \textbf{공동체의 실천(community
practice)}만이 규칙 적용에서 무엇이 옳고 그른지를 정립한다고
주장했다\autocite{kripke1982,rulefollowing2022}. 법에 대한 유비는 우리를
멈칫하게 한다. 만약 각 판사가 법률의 의미에 관하여 자신의 사적 판단만을
따른다면, 그 법률의 의미에 관한 객관적 사실은 전혀 존재하지 않게 될
것이며, 의미는 주관적 변덕으로 해체될 것이다. 오직 올바름에 대한 공적의
준거를 가진 \textbf{공동체의 실천}을 참조함으로써만, 우리는 한 판사의
해석이 \textsf{옳은지 그른지}를 말할 수 있다. 이러한 사유의 흐름은 해석적
수렴의 난제가 규칙 따르기 난제의 특수한 경우임을 시사한다. 해법은 개별
정신 내부에서 비밀스러운 지침을 찾는 데 있지 않고, 개인들이 \textsf{같은
길로 계속 나아가도록} 훈련하고 인도하는 사회적·제도적 틀을 살피는 데
있다.

이 지점에서 다음과 같은 이의를 제기할 수 있다. 어쩌면 법적 수렴은
환상이며, 우리는 법의 일관성을 과대평가하고 있는 것이 아닐까? 실제로
판사와 다른 공무담당자들이 종종 격렬히 의견을 달리한다는 것은 사실이
아닌가 ― 헌법재판에서 재판관들의 5--4 의견 분열이나, 법률 의미에 관한
법원 간 분열을 생각해 보라. 어쩌면 (비트겐슈타인의 회의론에 상응하는)
회의론자들이 옳을지도 모른다. 즉, \textsf{단 하나의 올바른 해석은 존재하지
않고}, 힘을 놓고 다투는 경쟁적 관점들만 있을 뿐이라는 것이다. 만약
그렇다면 수렴을 설명하는 과제는 무의미해질 것이다 ― 이를 설명할 수 있는
방법은 무정부 상태나 지속적 갈등뿐일 테니 말이다. 그러나 법체계가 실제로
기능하는 방식은 다른 그림을 제시한다. 주변적 쟁점에서의 불일치에도
불구하고, 법적 질문들의 압도적 다수에 관해서는 광범위한 합의가 존재하며,
불일치가 있는 곳에서도 제도적 절차들은 이를 해결하고 통일하려고
한다(상급심 심사, 전원합의체 심리, 법해석에 대한 입법적 보완 등).
판사들이 이유를 제시한 의견을 내고 출처를 인용한다는 사실 자체가, 한
해석을 더 낫거나 더 못하다고 판단할 수 있게 하는 공유된 논증 기준을
함의한다. 요컨대, 법적 실천의 참여자들은 올바르거나 적어도 \textsf{가장 잘
정당화된(most justified)} 해석이 존재하며 그것을 찾아낼 수 있다는
전제하에 행동한다. 그들은 보통 ``모두 주관적일 뿐''이라는 식으로
체념하지 않는다. 이러한 실천적 전제 ― 법적 질문들에는 이론상 추론적
방법론(reasoned methodology)을 통해 발견 가능한 \textsf{답(answers)}이
있다는 전제 ― 가 수렴의 가능성을 떠받친다. 우리는 극단적 형태의 해석적
회의론이 거짓(또는 최소한 전부는 아님)이라는 가정하에 나아갈 것이다. 곧,
법적 해석에는 대부분의 경우 수렴을 산출하는 유의미한 기준들이
\textsf{실제로} 존재하며, 그렇지 않은 경우에는 체계가 새로운 수렴을
확립하기 위해 작동한다는 것이다. 따라서 난제는 이러한 기준들과
메커니즘이 무엇인지를, 특히 수렴을 집행하면서도 역설적으로 이탈할 자유가
있는 듯 보이는 사법적 행위자의 문제적 역할을 고려하여, 분명히 서술하는
데 있다.

이제 \textbf{사법적 역할}에 관한 수수께끼를 명확히 해 보자. 사법적
행위자는 가장 대표되는 \textsf{해석자}, 개별 사건마다 법이 무엇인지
선언하는 제도적 목소리인 것 같다. 모든 일이 잘 흘러갈 때, 판사는 해석적
수렴을 결속시키는 접착제다. 판사는 모호성을 해소하고 일관된 적용을
보장한다. 그럼에도 판사는 잠재적 \textbf{해석적 분기}의 중심에 서기도
한다. 특이한 경향의 판사는 새로운 해석을 도입할 수 있으며, 그것이
제어되지 않는다면 그 해석이 전파될 수도 있다. 많은 법률가들과 학자들이
두려워하는 바는, 판사들이 해석의 외피를 두르고 사실상 법을 다시 쓸 수
있으며, 그리하여 수렴과 민주적 정당성을 모두 훼손할 수 있다는 점이다.
이러한 우려는 우리 이론에 대한 오해로 이어지기도 한다. 어떤 비평가들은
해석적 수렴 이론이 순진하게도 판사들이 \textsf{항상} 입법 의도에 따를
것이라고 가정한다고 생각할 수 있다. 이는 경험적으로 거짓인 듯하다.
반대로, 어떤 이들은 이 이론이 판사에게 너무 많은 재량을 부여하여,
``수렴''이란 결국 판사가 말하는 것이 곧 법이 되는 것에 불과하다고, 그
결과 이 이론이 공허하거나 순환적이 된다고 이해할 수도 있다. 이 두 가지는
실제로 결함이 될 수 있겠지만, 모두 이 이론에서의 사법적 역할을 오인한
데서 비롯된다. 난제를 해결하려면, 판사가 판단에 있어서는 독립적이면서도,
\textsf{그럼에도 불구하고} 어떻게 수렴의 방향으로 제약될 수 있는지를
설명해야 한다.

요약하면, 우리는 다면적 문제를 마주한다. (1) 법 내 의사소통은 언어의
일반적 모호성과 법 텍스트의 특유한 불확정성에 직면해 있다. (2) 법은
다양한 행위자들에 의해 구현되며, 이는 조정 문제를 야기한다. (3) 판사는
법을 해석하고 적응시키기 위한 자유가 필요하면서도, 동시에 공통의 방향을
고수해야 하는 딜레마를 예증한다. 다음 절은 분석적 해결을 전개함으로써
이러한 쟁점들을 정면으로 다룬다. 핵심적 전환은 의미의 \textsf{개인주의적}
도식(예: 입법자의 사적 의도, 판사의 개인적 철학)에서 의미의
\textsf{제도적}·\textsf{관행적} 도식(공동체적 실천과 역할로 정의된 의무를
통해 이해되는 규칙)으로 옮기는 것이다. 우리는 제대로 이해된 사법적
행위자가 규제받지 않는 의미의 창조자가 아니라, 공유된 기준을 통해 해석을
\textsf{규율하는(disciplines)} 실천의 중심적 참여자임을 보일 것이다.
단계적으로 이 이론을 정교화하여, 수렴이 어떻게 능동적으로 유지되는지를
보이고자 한다.

\section{단계적 해결: 개인적 의도에서 실천-매개적
화행적수인으로}

\subsection{1. 법에서의 의사소통: 화자-의미 모형을
넘어}

자연스러운 출발점은 법적 의사소통을 화자와 청자 간의 일반적 의사소통과
비교하는 것이다. 폴 그라이스(Paul Grice)의 의미론에 따르면, 성공적인
의사소통은 화자가 특정 의미를 전달하려는 의도를 가지며 청자가 그 의도를
인식하도록 의도하고, 청자가 실제로 그것을 인식할 때
발생한다\autocite{grice1957,grice1989}. 이를 유비하면, 입법부는
``화자''이고 그 의미가 법의 내용이며, 판사(또는 다른 공무담당자)들은 그
의도된 내용을 파악해야 하는 ``청자''라고 할 수 있다. 만약 모든 판사가
입법부의 의도를 올바르게 추론한다면, 해석적 수렴은 의도된 계획의
귀결로써 달성된다. 즉, 판사들의 화행적수인이 입법부의 의사소통적 의도와
일치하는 것이다.

그러나 이러한 순진한 형태의 화자-의미 모형(Speaker-Meaning Model)은 법의
경우에는 부적절하다. 일반적 대화와 달리, 법적 의사소통의 참여자들은
실시간으로 피드백을 주고받는 동등한 대화자가 아니다. 입법부는 보통
화자가 청자의 혼란을 즉각 바로잡을 수 있는 방식으로 판사의 오해를 직접
교정할 수 없다. 여기에 \textbf{시간적·맥락적 간극}이 존재한다. 법은 특정
시점에서 특정 집단에 의해 제정되지만, 나중에 전혀 다른 사람들이 초안
작성자가 전혀 보지 못한 상황들에 적용한다. 그라이스적 의미 기제 ― 화자의
의도와 청중의 그 의도의 인식 ― 는 이러한 조건 하에서 부담을 안는다.
첫째, 입법부는 단일한 정신(mind)을 가진 것이 아니며, 다수 구성원의
집단에 단일한 의도를 귀속시키는 것은 어려운 일이다(그것은 중앙값
투표자(median voter)의 의도인가, 모두가 수용할 수 있었던 타협인가,
추상적 ``객관적 의도''인가?). 둘째, 객관적 의도가 존재한다고
가정하더라도, \textbf{법의 청중은 분산되어 있으며(diffuse) 시간상 나중에
등장한다(deferred)}. 판사와 공무담당자들은 특정한 관행 ― 예를 들어,
``문언의 통상적 의미(plain meaning)'' 규칙, 법 해석의 구성 정범(canons
of construction) ― 을 따르는 전문적 해석 공동체에 속하며, 이 관행들이
그들이 의도를 추구하는 방식을 매개한다. 이는 곧 판사들이 인식하는 것이
원초적 심리적 의도가 아니라, 그러한 관행에 의해 걸러진 의도임을
의미한다. 본질적으로 의사소통의 연쇄는 \textbf{공적으로 접근 가능한
의미} ― 즉, 법의 텍스트와 그 텍스트를 해석하기 위한 법적 관행 ― 에 의해
매개된다.

이로써 우리는 단순한 그라이스적 구도에서 보다 \textbf{공적 의미론적
구도}로 이동한다. 즉, 법적 행위의 내용은 화자의 의지만으로 결정되지
않으며, 그 행위가 화행적수인되는 시점에서 공유된 언어적·법적 관행 안에서
\textbf{분별있게 이해될 수 있는} 방식에 의해 결정된다. 이 관점에서
의미는 \textsf{저작의 순간(at the moment of authorship)}에 고정되는
것이라기보다는, 어떤 특정 시점에서든 그 행위가 공동체의 공유된 해석 기준
속에서 화행적수인되고, 다투어지면서 안정화되는 것이다. 이는 익숙한 구별,
즉 \textbf{화자의 의미}와 \textbf{발화 (또는 문장) 의미} 사이의 차이를
반영한다\autocite{grice1989}. 입법기관은 특정한 목적(예: 대기오염 감소)을
갖고 그 목표를 달성하기 위한 문구를 선택할 수 있다. 그러나 그들이
제정하는 법률의 의미는, 상당 부분, 그 단어들의 관행적 의미론 및 법
공동체가 이해하는 법기술적 전문용어에 의해 고정된다. 사적 코드(private
code)가 아니라 공적으로 이해 가능한 방식으로 소통하기 위해, 입법기관이
그러한 관행들에 의도적으로 구속되려고 한다는 점은 전형적인 사실이다. 이
점에서 실패하는 경우, 예를 들어 불명확하거나 기묘한 언어를 사용하는
경우, 그 의도는 법으로 번역되지 못한다. G--S--F 틀 안에서 이러한 언어적
역량은 주로 \textbf{공유가능성} --- 즉, 공적 이해가능성과 공유된 어휘
내의 이유 제시 --- 에 속하며, 또한 규범 정식화 작업이 단발적 명령이 아닌
규칙을 구성하는 한에서는 \textbf{일반성}에도 관련된다. 극도로 모호하거나
상호 모순적인 법률은 정확히 이러한 이유로 인해 수렴을 저해하고
작동불능의 위험을 낳는다.

여기서 질문이 제기된다. 법의 의미가 본질적으로 공적인 것이라면, 입법
의도 자체를 논할 필요가 있을까? 일부 이론가들(소위
문언주의자(textualists))은 법 제정 당시의 \textsf{분별있는 독자가 이해하는
바}로서의 텍스트 의미에 초점을 맞추고, 주관적 의도는 배제해야 한다고
주장한다. 다른 이론가들은 텍스트가 모호할 경우 의도나 목적이 안내적
역할을 한다고 주장한다. 우리의 이론은 일정 부분 양쪽 입장을 수용할 수
있다. 즉, \textbf{해석의 출처들 간 위계}를 인정하는 것이다. 텍스트와
관행적 사용이 1차적 기준(공적으로 이용 가능한 내용)을 제공하고, 목적이나
의도에 관한 증거가 남아있는 모호성을 해결할 수 있다는 것이다. 수렴을
위해 핵심적인 점은 어떤 해석 전략이 사용되든, 그것이 \textbf{해석자들
사이에서 공유되어야 한다}는 것이다. 만약 한 체계 내의 절반의 판사들이
엄격한 문언주의자이고, 나머지 절반이 목적론자라면, 특정 사건에서 그들은
체계적으로 분기할 수 있다. 법체계는 일정 수준의 방법론적 다양성을 수용할
수 있지만, 안정적 수렴은 실제로는 지배적이거나 최소한 조정된 접근을
요구할 가능성이 높다. 실제로 수렴이 유지되는 방식 중 하나는
\textbf{방법론적 관행}이다. 판사들은 대체로 유사한 방식으로 법률을
해석하도록 훈련되며, 승인된 해석 정범과 선례를 사용한다. 이는 판사들의
화행적수인 함수들을 동조시키는(synchronize) 역할을 한다.

이 분석 단계에서 우리는 순수한 개인주의적 모형을 넘어 공적 의미와 공유된
해석 방법을 강조함으로써 개선을 이루었다. 그러나 우리는 아직 판사들이
수렴을 깨뜨리지 않고 재량을 어떻게 행사할 수 있는지를 완전히 해결하지
않았다. 이를 위해 우리는 이제 판사들의 해석을 둘러싼 \textbf{제도적
맥락과 피드백 메커니즘}을 고려해야 한다.

\subsection{2. 제도적 역할과 수렴의
규율}

모든 법적 행위자는 해석 산출물을 일정한 방향으로 유도하고 타자에 대한
책임성을 부과하는 \textbf{역할-구조화되고 실천-구조화하는
제약(role-structured and practice-structuring constraints)} 아래에서
활동한다. 여기서 ``제도적(institutional)''이라는 표현은 (``확립된 제도
내부''라는 의미가 아니라 이를 산출해 나가는) \textbf{과정적(processual)
의미}로 사용되고 있다. 즉, 이는 해석적 화행적수인이 시간의 경과에 따라
안정화되도록 만드는 침전된 역할, 관행, 그리고 피드백 메커니즘을 가리키는
것이지, 화행적수인에 선행하는 고정된 의미의 원천을 뜻하는 것이 아니다.
판사의 경우, 이러한 제약에는 해석 관행뿐 아니라 항소심 심사, 이유 제시
의무, 그리고 선례의 원칙이 포함된다. 이러한 요소들은 피드백 및 교정
메커니즘으로 기능한다. 예컨대 1심 판사의 특이한 해석은 항소심에서 파기될
수 있고, 최고법원에서도 판사는 다수의 지지를 확보하고 유지해야 하므로
합의적 규율이 작동한다. 또한 특정 해석의 흐름이 널리 공유된 규범적
기대와 현저히 어긋나거나 실질적으로 작동 불가능한 것으로 드러날 경우,
입법부는 개정이나 재입법을 통해 대응할 수 있다. 이처럼 법적 실천은 다른
행위자들에 의한 \textbf{화행적수인과 응답(uptake and response)}의
기회들로 관통되어 있으며, 이를 통해 경로 수정이 가능해진다. 중요한 점은,
여기서의 화행적수인은 단순한 찬동(endorsement)을 의미하지 않는다는
것이다. 그것은 따르기, 구별하기, 다투기, 좁히기, 교정하기 등, 공동체가
특정 행위의 의미를 공적으로 다루고 정합화하는 모든 방식을 포함한다.

이제 앞서 제시한 화행적수인 함수 \(U_A\)의 개념을 반복적으로 적용해
보자. 입법부 \(L\)이 법 \(x\)를 제정한다고 하자. 판사 \(J\)는 이를
\(y = U_J(x)\)로 해석하고, 사건에서 \(y\)를 적용한다. 이러한 사법적
적용은 다시 다른 행위자들의 화행적수인의 대상이 된다. 상급심은 이를
인용하거나 파기하거나 구별할 수 있고, 행정기관이나 검사들은 \(y\)에
맞추어 집행을 조정하거나 충돌을 야기하여 명확화를 촉발할 수 있다.
입법부가 \(y\)로의 해석의 전개에 강한 이의를 가지는 경우, 개정이나
재입법을 통해 \(x'\)를 제시하여 이후의 화행적수인을 사실상 \(z\)
방향으로 이동시키려 할 수 있다(예컨대, 지배적인 해석 방법 아래에서
\(U_J(x') \approx z\)가 되도록 설계하는 방식으로). 이러한 과정은
경제학적 의미의 단일한 ``평형상태(equilibrium)''에 반드시 도달하는 것은
아니다. 오히려 체계는 \textbf{국소적으로 안정된 초점점(locally stable
focal points)}으로 안정화되는 경향을 보인다. 즉, 현존하는 직업적 규범,
선례 관행, 심사 구조 아래에서 더 이상의 반복적 화행적수인과 피드백이
의미 있는 변화를 초래하지 않는 해석 패턴이 생산되는 것이다. 이러한
의미에서 해석적 수렴은 실천적 함의를 획득한다. 그것은 행위자 간, 그리고
시간의 경과에 따라 조정된 화행적수인이 미래의 행위를 인도하기에 충분히
견고한 공적 이해를 산출하여, 지속적인 재협상을 요구하지 않는 상태를
의미한다.

이제 법적 행위자의 \textbf{사회화와 훈련}을 고려해 보자. (Norman Malcolm
등의 해석에 따르면) 비트겐슈타인은 규칙 따르기가 오직 공동체 내의 훈련을
통해서만 가능하다고 강조했다. 그 공동체는 올바름(correctness)의 준거를
설정한다\autocite{malcolm1989}. 판사와 변호사는 법이론과 법실무에 대한 교육
및 연수과정 등에서 장기간의 훈련을 거치며, 이는 단순히 규칙을 배우는
것이 아니라 공동체 내 다른 이들처럼 \textsf{해석하고 추론하는 방법}을
배우는 것이다. 이는 우리가 \textbf{인식 공동체(epistemic community)}라고
부를 만한 것을 창출한다. 이 집단은 내부적으로 논쟁을 벌이기도 하지만,
무엇이 좋은 논증인지, 무엇이 개연성 있는 해석인지를 판단하는 기본적
전제를 공유한다. 이는 수렴에 결정적이다. 판사들이 의견을 달리할 때조차,
그들은 일정한 한계 내에 머물게 하는 법적 추론의 통용어(lingua franca)
내부에서 불일치를 드러내는 것이다. 예컨대, 판사가 개인적 사주나 별자리
등을 참조하여 법률을 해석하는 것은 사실상 상상하기 어려운 일이다. 이는
논리적으로 불가능해서가 아니라, 그러한 방법이 공동체가 공유하는 준거의
바깥에 있기 때문이다. 반대로, 입법 연혁을 참조하거나 선례에 따르는 것은
수용되는 해석 행위이다. 공동체는 정당성의 경계선을 설정함으로써
간접적으로 수렴을 강제한다. 즉, 그 경계를 넘어서는 해석은 전문적 비판,
판결의 파기, 또는 다른 사람들에 의한 준수 거부 등에 직면하게 된다. 이런
방식으로 \textbf{행위에서의 일치(agreement in acting)} ― 곧 ``어떻게
계속할 것인가''에 관한 공동체의 암묵적 합의 ― 가 법적 해석의 객관성을
뒷받침한다.

이제 이러한 논의들을 우리의 형식적 구성에 다시 연결해 보자. G--S--F
제약조건들은 우리가 \textbf{해석가능성 요건(interpretability
requirements)}이라 부를 수 있는 사항을 포함한다. 즉, 행위는 명료하고
이미 알려진 관행에 위치함으로써, \textbf{확정적인 화행적수인(determinate
uptake)}이 가능해야 한다. 입법적 행위의 경우, 이는 해당 텍스트가 공유된
언어와 해석 정범에 비추어 해당 법적 청중에게 이해 가능해야 한다는 것을
의미한다(\textbf{S}, 그리고 규칙적 형식(rule-like form)이 문제될 때는
\textbf{G}까지 포함). 사법적 행위의 경우, 실질적 성공의 한 조건은 그
결정이 법적 공동체의 다른 구성원들이 따를 수 있는 방식으로 정당화되는
것이다. 판결문의 의견들(opinions)은 단순히 결론을 기록하기 위해 작성되는
것이 아니라 \textbf{설득하고 설명하기 위해} 작성된다. 이와 같은 이유
제시를 결여한 판결은 여전히 사건을 종결시킬 수는 있지만,
\textbf{공유가능성(S)}을 충족하지 못하며 \textbf{초점점 안정성(F)}을
형성하지 못한다. 그 결과, 법적 행위자들은 그 판결을 어떻게 일반화하거나
존중해야 하는지 인식하지 못할 수 있다. 극단적 사안들에서, 설명되지
않았거나 불투명한 결정은 해석 협소화나 묵시적 불복을 초래하며, 그 권위를
약화시킨다.

수렴을 촉진하는 또 다른 제도적 요소는 \textbf{선례구속의
원칙}(\emph{stare decisis})이다. 이 원칙은 시간적 수렴을 직접적으로
촉진한다. 즉, 현재의 판사는 패턴을 깨뜨릴 강한 이유가 없는 한, 과거
판사들의 해석에 자신의 해석을 정렬시킨다. 선례는 화행적수인 함수의
기대가능한 연속성을 창출한다. 예를 들어 동일한 법률 \(x\)가 쟁점일 경우
특별한 변화가 없다면 2026년의 사법적 해석 \(U_{J_{2026}}(x)\)는 1996년의
사법적 해석 \(U_{J_{1996}}(x)\)와 대체로 동일해야 한다. 이러한 시간적
일관성은 예측 불가능성을 줄여 행위자 간 조정을 용이하게 한다. 각
행위자는 미래에 다른 이들이 해당 법률을 어떻게 해석할 것인지 예측할 수
있기 때문이다. 수렴의 관점에서 보면, 선례는 불일치의 차원들을 줄인다.
즉, 모든 판사가 법령을 매번 처음부터 재검토하여 광범위한 해석적 변동성이
발생할 수 있는 상황과 달리, 공통의 기준적 해석에서 출발하도록 만든다.

이러한 모든 메커니즘 --- 상급심에 의한 심사, 입법적 파기, 전문적 훈련,
이유가 제시된 판결의 의견들, 그리고 선례 --- 은 논문의 제목에서 언급한
\textbf{법적 실천의 구조(structure of legal practice)}를 구성한다.
주의할 것은, 이 `구조' 자체는 과거의 행위와 화행적수인이 축적된
침전물이라는 점, 즉 위로부터 부과된 선행하는 틀이 아니라 의사소통적
교정이 제도화된 결과물이라는 점이다. 이들은 행위자들 간, 그리고 시간의
경과에 따라 해석적 산출물이 정렬되도록 집합적으로 기능한다. 유용한 비유
하나는 \textbf{오류 교정(error-correction)}이다. 개별 행위자가
이탈하더라도, 구조 자체가 그 오류를 교정하거나 억제하여 시스템 전체가
궤도에서 벗어나지 않도록 유지한다. 해석적 수렴 이론은 이러한 구조적
특성들이 우연적 요소가 아니라 근본적이라고 주장한다. 바로 이러한
구조들이 법적 실천을 협력적 기획으로서 가능하게 만드는 것이다.

\subsection{3. 사법적 행위자의 역할 명료화: 해석에서의 권위와
의무}

이러한 실천-구조화된 관점에서 볼 때, 자주 오해되어 온 사법적 행위자의
역할을 명료화할 수 있다. 판사는 흔히 두 가지 상반된 방식으로 묘사된다.
하나는 판사를 입법자의 의도나 제정자의 의도를 충실히 구현해야 하는
\textsf{종속적 행위자}, 즉 대변인으로 보는 견해이고, 다른 하나는 판사를
자신의 판결을 통해 법의 의미를 일방적으로 고정하는 \textsf{독립적 권위자},
즉 입법자로 보는 견해이다. 여기에서 전개된 해석적 수렴 이론은 판사를
\textbf{매개적 행위자}로 재개념화함으로써 이 두 극단을 해소한다. 판사는
법전상의 법과 실제로 작동하는 법 사이를 매개하며, 그 최상위 의무는 법을
타자들이 신뢰할 수 있게 수용할 수 있는 정합적이고 공적으로 이해 가능한
체계로 유지하는 데 있다.

이 견해에서 사법적 권위는 법적 실천 전반에 대해 형이상학적
우선성(metaphysical priority)을 갖는 것이 아니라 \textbf{획득된
권위(earned authority)}의 한 형태라는 점이 중요하다. 그것은 오직
판결들이 공동체의 수렴적 실천 속에서 초점점들(focal points)로
화행적수인되는 한에서만 확장된다. 따라서 법원은 그 산출물이 후속하는
화행적수인을 통해서만 --- 즉, 전문가 공동체가 적절하다고 간주하는
방식으로 그 판결이 준수되고, 일반화되며, 구별되고, 범위가 제한되거나,
비판되거나, 공고화됨을 통해서만 --- 법형성적(law-shaping) 지위를 획득할
수 있는 분쟁 해결 제도로 가장 잘 이해된다. 이에 따라 \textbf{판사를
직무-구속적 해석자(office-bound interpreter)}로 말하는 것이 보다
정확하다. 판사의 자유는 사적 개인의 자유가 아니라,
불편부당성(impartiality), 이유제시(reason-giving), 승인된 법적 원천에
대한 충실성(fidelity to recognized sources), 심사에 대한
응답성(responsiveness to review)과 같은 \textsf{역할-책임들에 의해 제약된}
직무의 자유이다. 이러한 책임들은 재판을 기계적인 것으로 만들지 않는다.
판단(judgment)은 불가결한 것으로 남아 있지만, 순전히 기묘하거나, 또는
결론을 정해두고 추론을 조작하는 식의 일탈은 배제한다. 우리의 틀에서
판사의 화행적수인 함수 \(U_J\)는 백지위임장이 아니다, 그것은 훈련,
공유된 해석 규범, 그리고 위에서 서술한 감시 메커니즘에 의해
조건지어지며, 이 모든 것은 해석적 수렴의 유지를 (그리고 필요할 경우
책임있는 전환을) 향하고 있다.

사법적 역할에 관한 오해는 흔히 이 스펙트럼의 한쪽 극단에만 집중함으로써
발생한다. 비판자들은 판사에게 창조의 여지가 조금이라도 주어진다면 수렴은
버려지고, 법은 단순히 사건마다 판사가 만든 ``판사-제정법(judge-made
law)''이 된다고 우려한다. 그러나 우리가 제시하는 이론은 그 두려움이
과장되었음을 보여준다. 창조의 여지는 수렴적 실천에의 전념(commitment)에
의해 제한된다. 판사들은 자신의 판결이 다른 판결들과 조화를 이루고, 다른
이들이 따르기를 \textsf{원한다(want)}. 이는 직업적 에토스에서 나오기도
하고, 효과없는 판결로 달성되는 것이 거의 없기 때문이기도 하다. 경제적
용어로 말하면, 판사는 해석의 많은 외부 효과를 내부화한다. 즉, 자신이
지나치게 일탈하면 상급심 법원이나 입법부가 반응할 것이고, 자신의 명성이
손상될 것임을 알고 있다. 반대로, 어떤 이들은 판사가 지나치게 구속되어
(즉 단순히 충실한 대리인이 되어) 새로운 상황에 대응하거나 부정의를
교정하지 못해, 법이 경직되고 반응하지 못할 것을 우려한다. 바로 이
지점에서 판사의 이유에 기초한 재량(reasoned discretion)의 역할이
나타난다. 판사는 법이 개연성 있는 해석의 범위 내에 적응하도록 밀어붙일
수 있다. 만약 그들의 적응적 해석이 사회적 가치나 법적 원칙과 진정으로
공명한다면, 그것은 (미래의 법원, 여론, 심지어 입법부의 묵시적 승인과
같은) \textsf{다른 행위자들에 의해 수용될} 가능성이 크다. 그 경우 수렴은
새로운 평형상태로 이동한다 ― 비록 그것이 원래 입법부가 명시적으로 상상한
바는 아니더라도, 법체계 전반적 차원의 전념들과 일관된 상태로 남아 있게
된다.

구체적 사례가 도움이 될 수 있을 것이다. 적용범위가 넓은 포괄적 표현으로
수십 년 전에 제정된 법률이, 문자 그대로 해석할 경우 변화된 상황 속에서
터무니없거나 부당한 결과를 초래하는 경우를 생각해 보자. 판사는 딜레마에
직면한다. 법문의 글자에 충실히 머물러 부당한 결과를 산출할 것인가(충실한
대리인), 아니면 불합리한 결과를 피하기 위해 해석을 조정할 것인가(창의적
입법자)? 이 지점에서 다른 판사들은 서로 다른 결정을 내릴 수 있으며, 바로
여기서 해석적 분기가 위협으로 다가온다. 그러나 제도화된 틀이 해결을
향도한다는 점에 주목할 필요가 있다. 예컨대 판사가 입법부가 의도하지
않았을 결과를 피하기 위해 문자적 텍스트에서 벗어나는 것을 허용하는
\textbf{부당한 결과 회피 원칙(absurdity doctrine)}이 있다. 이는 사실상
더 높은 수준의 입법 의도 ― 즉, 부당한 결과를 초래하지 않으려는 의도 ― 를
예우할 것을 요청하는 것이다. 또한, 헌법적 가치나 적용가능한 국제 규범을
고려하여 법률을 해석하는 기법도 있다. 만약 판사 A는 부당한 결과를 피하기
위해 문언에서 일탈하는데 판사 B는 엄격히 문언을 고수한다면, 항소심
법원은 충돌을 해소하기 위해 개입하고자 할 수 있으며 그 중 하나의 방식을
정당한 것으로 승인할 수 있다. 만약 항소심 법원이 부당한 결과 회피를
받아들인다면, 그 해석을 수렴적 실천에 연결하는 논거도 제시하려고 할
것이다. 예를 들어, ``어느 분별있는 입법자도 이러한 결과를 의도하지
않았을 것이므로, 비록 우리의 해석이 가장 문언적이지는 않을 지라도, 해당
법률의 목적과 법체계의 합리성을 더 잘 예우한다''는 식이다. 반대로 법원이
문언을 고수한다면, 이는 입법부에 다음과 같은 신호를 보내는 것일 수 있다.
``우리는 판사로서 법문언에 구속을 받는다. 만약 그 구속의 결과가
나쁘다면, 당신(입법부)이 법을 변경하는 수밖에 없다.'' 어느 경우든
대화(dialogue)가 존재한다. 판사와 입법부는 상호 조정을 통해 궁극적으로
일관된 법 질서를 유지한다. 수렴은 정적인 동일성을 의미하지 않는다.
그것은 조정된(coordinated) 진화의 과정을 포함할 수 있다.

이러한 설명에서 도출되는 중요한 귀결은, 사법적 행위자의 역할이
본질적으로 \textbf{중간적(interstitial)}이며 동시에
\textbf{상호작용적(interactive)}이라는 점이다. 판사들은 법전에 존재하는
법과 재판이라는 실천적 결과 사이에 위치하며, 그들의 충실성은 이중적이다.
즉, 한편으로는 법률과 선례와 같은 법의 인정된 원천(recognized sources of
law)에 대한 것이며, 다른 한편으로는 정의와 이성에 대한 법의 열망에 대한
것이다. 이러한 이중적 충실성이야말로 총체적 해석적 분기나 형식주의의
무비판적 수용 어느 쪽으로도 흐르지 않도록 막아 준다. 어떤 의미에서
판사들은 번역자와 같다. 즉, 그들은 법적으로 이용 가능한 자료들---법률,
판례, 그리고 기타 인정된 원천들---을 구체적 사건의 맥락 속으로 번역한다.
좋은 번역은 원문 텍스트의 의미를 충실히 유지함과 동시에, 목표
언어(target language) 내에서 사리에 맞게 만드는 것을
요구한다\footnote{여기서 목표 ``언어''란 구체적 사실관계와 사회적 맥락에
  해당한다.}. 번역이 목표 청중에 대해 받아들일 만하지 못하면 실패이며,
원문을 배반해도 실패다. 판사들은 두 가지 실패 모두를 피하려 애쓴다. 가장
뛰어난 판사들은 법의 문언과 목적을 존중하면서도 맥락 속에서 설득력 있게
들리는 판결을 내린다. 바로 그러한 판결이 다른 행위자들로부터 수렴을
이끌어내며, 모두가 ``그렇다, 이것이 법이 되어야 할 모습으로
\textsf{받아들일 만하다}''고 말할 수 있게 된다.

이제 우리는 사법적 독립과 해석적 수렴이 대립한다는 \textsf{오해}를 불식할
수 있다. 오히려, 올바로 이해된 사법적 독립은 수렴의 구성요소다. 정치적
압력이나 개인적 이해관계와 같은 부당한 영향으로부터 독립된 판사는, 다른
이들이 정당하다고 인정하는 \textsf{법적} 이유와 공유된 해석 원칙에 더
충실할 가능성이 크다. 판사가 불편부당하고 정합적일 때, 기묘하거나 당파적
편향없이 공통된 법적 이유에 기초한 해석을 통해 수렴을 강화하게 된다.
만약 판사들이 선거나 개인적 고려에 따라 결과를 달리한다면 공정성이
위협받고 법의 일관성이 훼손될 것이다. 따라서 사법부를 특정한
압력으로부터 격리하는 것은, 해석적 실천이 법적 이유에 따른 타당성(legal
merits)에 기초한 수렴 상태를 유지하게끔 하는 것을 목표로
한다\autocite{ross2013reconsidering}.

마지막으로 극단적 경우를 생각해 보자. 만약 사법부 전체가 법체계(또는
사회)가 수용할 수 없는 방향으로 일탈하기 시작한다면 어떻게 될까? 예를
들어, 법원이 어떤 법률에 대해 그 명백한 목적을 무력화하는 매우 기이한
해석을 일률적으로 채택한다고 가정하자. 이때 최종적인 수렴을 뒷받침하는
안전장치는 \textbf{입법적 교정}이다. 입헌적 민주주의 체제에서는
입법부(또는 헌법 개정을 통한 국민)가 법규범에 대한 마지막 발언자라고도
할 수 있다. 입법부는 법을 명확히 하거나 다시 작성하여 원래 의도된 해석을
재정립할 수 있다. 이 권한이 존재한다는 사실은 판사들에게 미묘하지만
실제적인 영향을 미친다. 즉, 판사들은 심각한 오독이 입법에 의해 뒤집힐 수
있고(그리고 사법부를 곤혹스럽게 할 수 있음)을 알기에, 개연성 있는 범위
내에 머무를 이유가 있다. 반대로, 정치적 기관들이 사법부의 해석을
묵인한다면, 그 묵인은 일종의 인가로 기능하여, 사법부의 해석을 실효적
법(effective law)에 통합한다. 시간의 경과에 따라 이는 모두가 그 해석에
맞춰 이해를 조정하는 수렴적 실천의 일부가 된다.

이 단계별 해결의 결론은 다음과 같다. 고립된 의도적 또는 의미적
행위로부터 화행적수인의 실천-매개적 구조로 우리의 관점을 일단 옮기게
되면, 해석적 무정부 상태의 우려는 대체로 사라지게 된다. 사법적 해석은 더
이상 통제불가능한 분기의 원천으로 보이지 않고, 오히려 입법부, 법원,
행정부 공무담당자, 그리고 법의 수범자들이 서로의 해석을 상호적으로
형성하고, 다투고, 안정화하는 더 넓은 의사소통적 과정의 한 국면으로
드러난다. 따라서 사법적 행위자는 법의 최종적 의미에 대한 형이상학적
판정자(metaphysical arbiter)가 아니라, 법적 실천을 공적으로 이해가능하고
규범적으로 일관적인 것으로 유지할 책무가 부여된 직무-구속적
매개자(office-bound mediator)로 나타난다. 해석적 수렴은 행위 단계에서의
G-S-F 적정성 조건 제약과 더불어, 공유된 사회적·방법론적 관습, 그리고
제도적으로 구조화된 화행적수인, 응답, 심사의 형식들이 상호작용함으로써
유지된다. 따라서 일탈은 불가능한 것도 아니고 치명적인 것도 아니다. 여러
일탈들은 이유 제시, 선례, 항소심에 의한 감독, 입법적 응답, 그리고
전문적인 학술적 비평을 통하여 다투어질 수 있고, 교정될 수 있으며, 대개
수렴 안에 흡수가능한 상태가 된다. 그 결과는 정태적인 만장일치의 모습이
아니라, 행위자들 사이의 조정된 화행적수인을 통해 법적 의미가 시간의
경과에 따라 안정화되는 \textbf{자기-교정적 의사소통 질서}의 모습이다.
수렴의 수수께끼는 하나의 단일한 주권적 해석자를 찾아냄으로써 해결되는
것이 아니라, 그 자체가 과거의 행위들과 화행적 수인의 침전물인 법적
실천의 구조가, 공적으로 공유 가능한 법적 이해의 초점점들을 지속적으로
생성하고 유지하는 방식을 보여줌으로써 해결된다.

\section{결론}

이 글은 다음과 같은 하나의 난제에서 출발하였다. ``하나의 법적 질서가
서로 다른 역할을 수행하는 복수의 행위자들에 의해 산출되고, 해석되며,
적용됨에도 불구하고, 어떻게 그 법질서는 공적으로 이해 가능하고
규범적으로 정합적인 상태를 유지할 수 있는가?'' 내가 주장한 바에 따르면,
그 답변은 어떤 단일한 주권적 화자의 의미가 체계 전체를 통해 단순히
전달된다는 식의 허구에도 있지 않고, 제약 없는 해석적 다원주의에 대한
회의주의적 그림에도 있지 않다. 법은 공유된 해석 기준 안에서 법적
행위들이 화행적 수인되고, 다투어지고, 일반화되며, 신뢰의 대상으로 기능할
수 있는 한에서만 법적 힘을 획득하고 유지하게 되는, 실천-매개적인
화행적수인의 과정을 통해 결속된다.

이 출발점으로부터, 이 글은 두 가지 주장을 전개하였다. 첫째, 법적 행위는
일반성, 공유가능성, 초점점 안정성이라는 적정성 제약을 충족하는 한에서만
법이 되는 의사소통적 작용으로 이해되는 것이 가장 적절하다. 이러한 제약은
왜 일부 제정행위, 판결, 행정적 조치, 그리고 준수의 양식들은 작동하는
규범으로 안정화되는 반면, 다른 것들은 주변적인 위치에 머물거나
소멸하는지를 설명한다. 둘째, 법적 행위자들은 ``형성''과 ``해석''이라는
고립된 영역으로 분할되지 않는다. 입법부, 행정부 공무담당자, 법원, 그리고
법의 수범자들은 각기 상이하고 불균등한 방식으로, 자신들의 고유한
화행적수인 형식을 통하여 계속적인 법형성 과정에 참여한다. 따라서
권력분립은 법형성과 법해석 사이의 형이상학적 분할이라기보다, 공통의 실천
내부에서 역할, 책임, 상호 교정이 분화되어 있는 구조로 이해된다.

이와 같은 재구성은 사법적 역할을 명확하게 해 준다. 법원은 법의 최종적
의미에 대한 형이상학적 신탁자가 아니라, 법 공동체 내부에서의 후속적
화행적 수인을 통해서만 그 결정이 법형성적 초점점이 될 수 있는 분쟁-해결
제도다. 따라서 사법적 권위는 자기정초적인 것이 아니라, 매개되고 획득되는
것이다. 올바르게 이해된 사법적 독립은 수렴에 대한 위협이 아니라 오히려
그 조건들 중 하나이다. 왜냐하면 그것은 사법적 이유가 당파적 의지나 사적
선호가 아니라 공통의 법적 기준에 대한 답변가능성을 유지하도록 해 주기
때문이다. 이와 동시에, 이유 제시, 선례, 항소심에 의한 감독, 입법적 응답,
그리고 전문적인 학술적 비평은 사법적 혁신이 분열로 해체되는 것을 막는다.

따라서 전체적인 모습은 정태적 만장일치의 모습이 아니라 스스로 교정해
나가는 의사소통적 질서의 모습이다. 법적 실천의 구조 자체는 과거의
행위들과 화행적 수인의 침전물이지, 위로부터 부과되는 선재하는 틀이
아니다. 법적 의미가 공통적인 것으로 남아 있는 것은 모든 참여자들이
동일하게 사고하기 때문이 아니라, 불일치가 공적으로 공유 가능한 형식들
속으로 유도되어 비판되고, 수정되며, 이해의 초점점들로 안정화될 수 있기
때문이다. 이러한 관점에서 해석적 수렴은 법리학의 신비로운 배경적 가정이
아니라, 법이 이를 통해 공통의 규범적 실천이 되고 또 그러한 것으로
유지되게 하는 지속적인 성취이다.

\clearpage

\printbibliography

\clearpage

\begin{abstract}

This article develops a second step in an interpretive-convergence
approach to jurisprudence by shifting the focus from the existence
conditions of law to the \textbf{structure of legal practice} and the
roles of different legal agents. Building on an earlier account that
understands legal normativity as arising from \textbf{interpretive
convergence} within a community, this paper analyzes how concrete
\textbf{legal acts} and institutional \textbf{offices} contribute to
that convergence, and how the judiciary's role has been systematically
misunderstood.

The first part develops a general theory of \textbf{legal acts}. A legal
act is understood as a communicative operation through which law is
brought into being and counts as law only insofar as it satisfies three
felicity constraints: \textbf{Generality}, \textbf{Shareability}, and
\textbf{Focal-point Stability} (the G--S--F conditions). These
constraints explain why certain legislative enactments, judicial
decisions, administrative measures, and practices of compliance
stabilize as legal norms, while others remain peripheral or gradually
fade. The second part reconceptualizes the principal \textbf{legal
agents}---legislatures, executive officials, courts, and legal subjects
or citizens---as participants in an ongoing process of law-formation
insofar as they contribute to interpretive convergence. In doing so, it
challenges the orthodox division according to which ``making'' is the
function of legislatures and ``interpreting'' the function of courts.

The core of the paper addresses the \textbf{judicial role}. It argues
that courts are best understood not as metaphysical oracles of law's
final meaning, but as \textbf{dispute-resolution institutions} whose
decisions acquire broader authority only through subsequent uptake
within the legal community. Apparent puzzles about judicial supremacy,
hard cases, and ``activism'' are reframed as questions about how far
particular decisions succeed in meeting the G--S--F constraints under
existing professional norms, precedential practices, and structures of
review. The final sections draw out implications for separation of
powers, the discipline of precedent, and the legitimacy of judicial
reasoning. Rather than re-stating the foundations of interpretive
convergence, the article shows how that framework reconfigures our
understanding of \textbf{legal acts, legal agents, and the internal
architecture of legal practice}.

\end{abstract}

\end{document}
